\documentclass[11pt]{article}
\usepackage[T1]{fontenc}
\usepackage{lmodern,amsmath,amssymb,amsthm,microtype}
\usepackage[margin=1in]{geometry}
\usepackage[colorlinks=true,allcolors=blue]{hyperref}
\hypersetup{pdftitle={Negative answers to the printed totally mixed face questions},pdfauthor={Alec Kriebel}}
\newtheorem{theorem}{Theorem}
\newcommand{\R}{\mathbb R}
\newcommand{\Q}{\mathbb Q}
\newcommand{\conv}{\operatorname{conv}}
\newcommand{\dimn}{\operatorname{dim}}
\title{Negative answers to the printed\\totally mixed face questions}
\author{Alec Kriebel\thanks{Independent researcher. ORCID:
\href{https://orcid.org/0009-0001-9320-500X}{0009-0001-9320-500X}.}}
\date{October 17, 2026 \quad Version 1.0.0}
\begin{document}
\maketitle
\begin{abstract}
For two polytopes, retain the faces of their Minkowski sum whose two
canonical summands have positive dimensions and whose dimension is the
sum of those dimensions. This is the definition in Payne's printed
Oberwolfach question. Let $D$ be the standard four-simplex and $I=[0,1]$.
For the full-dimensional lattice polytopes $D\times I$ and $(-D)\times I$
in $\R^5$, the qualifying-face order complex is a disjoint union of two
connected pure graphs, each with $50$ vertices and $60$ edges. Its
rational Betti numbers are $(2,22,0,\ldots)$. Thus both printed
assertions, shellability and the homology of a wedge of three-spheres,
fail. The shared interval factor deliberately produces a nongeneric
example. The argument concerns the printed additive-dimension
definition; it does not settle shellability for the different mixed
faces used in stable tropical intersections.
\end{abstract}

\section{The exact question and the example}
Payne's question \cite[p.~2464, Problem~3]{P} concerns polytopes
$P_1,P_2\subset\R^n$ and those faces $F=F_1+F_2$ satisfying
\begin{equation}\label{filter}
 \dimn F_i>0\quad(i=1,2),\qquad
 \dimn F=\dimn F_1+\dimn F_2 .
\end{equation}
It asks whether their inclusion poset is shellable and whether its
realization has the rational homology of a wedge of $(n-2)$-spheres.
There is no stated general-position assumption. We use the standard
order complex: vertices are qualifying faces and simplices are
nonempty strict chains. No artificial extrema are added.

For a face exposed by a functional $u$, its canonical summands are the
faces $P_i^u$ exposed by the same $u$. The standard uniqueness theorem
for Minkowski-sum faces appears, for example, in
\cite[Lemma~4 of the cited arXiv version]{FW}. It also removes an existential ambiguity in
\eqref{filter}: if nonempty summand faces $G_i$ have
$G_1+G_2=P_1^u+P_2^u$, exposure first forces $G_i\subseteq P_i^u$.
For every other functional, the two nonnegative deficits between
their support functions sum to zero. Each deficit is therefore zero,
so $G_i=P_i^u$.

\begin{theorem}\label{main}
Let $D=\conv(0,e_1,e_2,e_3,e_4)\subset\R^4$ and $I=[0,1]$.
For $P_1=D\times I$ and $P_2=(-D)\times I$, let $T$ be the
inclusion poset selected by \eqref{filter}. Its order complex is a
disjoint union $K\sqcup K$, where $K$ is a connected pure graph with
$50$ vertices and $60$ edges. In particular,
\[
 H_0(|T|;\Q)=\Q^2,\qquad H_1(|T|;\Q)=\Q^{22},\qquad
 H_j(|T|;\Q)=0\quad(j\ge2).
\]
The poset is not shellable, under either pure or standard nonpure
shellability, and its realization does not have the rational homology
asked for in the printed question.
\end{theorem}

Both summands are full-dimensional lattice polytopes in $\R^5$.
The proof uses neither empty realizations nor low-dimensional sphere
conventions. The simplex-difference geometry is classical
\cite{FW,ST}; the observation used here is the effect of a common
interval factor on the exact filter \eqref{filter}.

\section{Complete face calculation}
Write $v_0=0$, $v_i=e_i$, and $E=\{0,1,2,3,4\}$.
Every nonempty proper face of $D-D$ is uniquely of the form
\[
 F(A,B)=\conv\{v_i:i\in A\}+\conv\{-v_j:j\in B\},
 \qquad \varnothing\ne A,B\subset E,\quad A\cap B=\varnothing .
\]
Indeed, for a nonzero functional $u$, take $A$ and $B$ to be its
maximum and minimum index sets. They are disjoint because the vertices
affinely span $\R^4$. Conversely, assign values $1,-1,0$ on $A,B$ and
the remainder, and subtract the value at $v_0$ to obtain a linear
functional with exactly these extrema.

The two summand direction spaces intersect only at zero: a vector
in their intersection gives an affine dependence on the disjoint
vertex sets, with coefficient sum zero on each set. Affine
independence makes every coefficient zero. Consequently,
\begin{equation}\label{dim}
 \dimn F(A,B)=|A|+|B|-2,\qquad
 \dimn D^u=|A|-1,\quad \dimn(-D)^u=|B|-1 .
\end{equation}
Moreover,
\begin{equation}\label{order}
 F(A,B)\subseteq F(C,L)\quad\Longleftrightarrow\quad
 A\subseteq C,\ B\subseteq L .
\end{equation}
For the nontrivial direction, evaluate every $v_i-v_j$,
$i\in A,j\in B$, against a functional exposing $F(C,L)$.
Attaining the maximum difference forces $i\in C,j\in L$.
The reverse direction follows directly from convex hulls.

Thus the base qualifiers have sizes $(|A|,|B|)=(2,2),(2,3),(3,2)$.
There are $\binom52\binom32=30$ minimal elements of type $(2,2)$
and $2\binom52=20$ maximal elements. Each maximum has three lower
minima, and each minimum has two upper maxima, obtained by placing
its one unused index on either side. There are no three-element
strict chains. The base order complex $K$ is therefore a pure graph
on $50$ vertices with $60$ edges.

To check connectedness, suppress the degree-two minimal vertices.
The resulting graph has signed two-subsets $A^+,A^-$ of $E$ as
vertices, with $A^+$ adjacent to $B^-$ exactly when $A\cap B$ is empty.
Two distinct intersecting same-sign subsets have the opposite-sign
complement of their union as a common neighbor. If two same-sign
subsets are disjoint, a two-subset containing one element of each
intersects both; apply the preceding path twice. All vertices of a
fixed sign therefore connect, and every vertex has a neighbor of
the other sign. The graph, and hence its subdivision $K$, is connected.
Contracting a spanning tree shows that $K$ is homotopy equivalent
to a wedge of $60-50+1=11$ circles.

Now $P_1+P_2=(D-D)\times[0,2]$. A functional $(u,t)$ with $t>0$
or $t<0$ exposes an upper or lower endpoint face; its summand
dimensions and qualification are exactly the base ones.
If $t=0$, write $a,b$ for the base summand dimensions and $c$ for
the base sum dimension. The two summands then contain the interval
factor, so their dimensions sum to $a+b+2$, while their sum has
dimension $c+1\le a+b+1$. No such face passes \eqref{filter}.
If $u=0,t\ne0$, the additional horizontal end face has dimension
$4$ but both summands have dimension $4$, so it also fails.
The whole sum has dimensions $(5,5,5)$ and fails; the empty face
fails positivity.

These cases exhaust all faces. The qualifying poset is precisely
two endpoint copies of the base poset, with no cross-comparisons.
This proves the asserted decomposition and homology.

In any proposed shelling, the first edge from the component different
from the initial edge has empty intersection with the preceding
facets. That intersection has dimension $-1$, whereas a new edge
requires a pure zero-dimensional intersection
\cite[Definition~2.1]{BW}. This rules out both shellability conventions.
Finally, a wedge of three-spheres, including the zero-fold point,
is connected and has no first homology. The same obstruction would
persist if one replaced the printed sphere dimension by $n-3$.

\section{Related results, scope, and verification}
For comparison, $D_3$ and $-D_3$ in $\R^3$ yield just six incomparable
qualifying faces: the index pairs partition four vertices into two
pairs. Their order complex is six points. This refutes the printed
wedge-of-circles homology assertion but is shellable in dimension zero.
Scaled opposite tetrahedra and their six-ray geometry already occur
in Sturmfels--Tevelev \cite[Examples~1.3 and 4.10]{ST}; this lower-dimensional control is not presented as new.
The five-dimensional example is needed here for nonshellability.

The term mixed has several different definitions. In the published
Adiprasito et al.\ paper \cite[Section~1.2]{AB}, totally mixed faces
require facets of the summands. That counting theorem is not used.
The public 2020 manuscript underlying Adiprasito's 2024 article
\cite[p.~2, Corollary~1.3 of the public manuscript]{A} instead selects
faces satisfying
\[
 \dimn\Bigl(\sum_{i\in J}F_i\Bigr)\ge |J|
 \quad\text{for every nonempty index subset }J.
\]
For two summands this is broader than \eqref{filter}.
For instance, a proper vertical face $F(A,B)\times[0,2]$ with
$|A|=|B|=2$ has summand dimensions $2,2$ and sum dimension $3$.
It meets that lower-bound criterion but fails additivity. Such
vertical faces are absent from our poset. The counterexample therefore
does not contradict the different stable-intersection theorem, or
resolve its shellability question.
Related tropical link obstructions occur in Hiatt--Simpson--Wang--Wu's
revised 2026 construction \cite[Section~8]{HSWW}. Its Sch\"on variety
in $(\mathbb C^*)^8$ is a different object; no identification with
this two-summand exact face poset is supplied or assumed here.

The supplement records a bounded primary-source priority audit,
including current versions and related tropical results. No earlier counterexample or theorem settling nonshellability for
this exact poset was identified. Such a result, together with the
classical homology antecedent, would already settle both assertions. This is not an exhaustive
literature or first-ever priority certificate. The 2024 journal
metadata for \cite{A} was verified, but its paywalled final body was
not inspected; the definition comparison above is explicitly with
the public 2020 manuscript.

Two independently authored standard-library programs accompany the
analytic proof. One reconstructs facets and their intersection lattice
from integer coordinates, all $543$ nonempty sum faces, and all $8649$
summand-face pairs. The other independently constructs the order complex
and its signed rational boundaries, obtaining boundary rank $98$.
Both run with explicit exceptions in ordinary and optimized Python;
positive and adverse controls are included. These computations are
supplementary, and the proof above requires no program.

\textbf{AI and review disclosure.} AI tools were used extensively in
developing, verifying, researching, and writing this result and its
supporting materials. This preprint is unrefereed. Independent
adversarial AI reviews do not constitute external human peer review,
journal acceptance, or formal proof-assistant certification.

\clearpage
\begin{thebibliography}{9}
\bibitem{P} S.~Payne, \emph{The complex of totally mixed faces},
Problem~3 in J.~Doolittle (collector), Open Problems, Oberwolfach
Report~39/2019, p.~2464.
\href{https://doi.org/10.4171/OWR/2019/39}{doi:10.4171/OWR/2019/39}.
\bibitem{FW} K.~Fukuda and C.~Weibel, \emph{On $f$-vectors of Minkowski
additions of convex polytopes}, Discrete Comput. Geom.\
\textbf{37} (2007), 503--516.
\href{https://arxiv.org/abs/math/0510470v4}{arXiv:math/0510470v4}.
\bibitem{ST} B.~Sturmfels and J.~Tevelev, \emph{Elimination theory for
tropical varieties}, Math. Res. Lett.\ \textbf{15} (2008), 543--562.
\href{https://arxiv.org/abs/0704.3471}{arXiv:0704.3471}.
\bibitem{BW} A.~Bj\"orner and M.~L.~Wachs, \emph{Shellable nonpure
complexes and posets. I}, Trans. Amer. Math. Soc.\
\textbf{348} (1996), 1299--1327.
\href{https://doi.org/10.1090/S0002-9947-96-01534-6}{doi:10.1090/S0002-9947-96-01534-6}.
\bibitem{AB} K.~Adiprasito, P.~Brinkmann, A.~Padrol, P.~Pat\'ak,
Z.~Pat\'akov\'a, and R.~Sanyal, \emph{Colorful simplicial depth,
Minkowski sums, and generalized Gale transforms}, Int. Math. Res.
Not.\ \textbf{2019}(6), 1894--1919 (online 2017).
\href{https://doi.org/10.1093/imrn/rnx184}{doi:10.1093/imrn/rnx184}.
\bibitem{A} K.~Adiprasito, \emph{A note on the topology of Minkowski
sums and complete intersections}, Pure Appl. Math. Q.\
\textbf{20}(1) (2024), 17--28.
\href{https://doi.org/10.4310/PAMQ.2024.v20.n1.a2}{doi:10.4310/PAMQ.2024.v20.n1.a2}.
Inspected public manuscript: \emph{A note on the topology of polyhedral
hypersurfaces and complete intersections}, July~27, 2020,
\href{https://webusers.imj-prg.fr/~karim.adiprasito/cint.pdf}{author PDF}.
\bibitem{HSWW} S.~Hiatt, C.~Simpson, B.~Wang, and C.~Wu,
\emph{Topology and Euler characteristics of tropical varieties},
arXiv:2606.10817v2, July~24, 2026.
\href{https://arxiv.org/abs/2606.10817v2}{arXiv:2606.10817v2}.
\end{thebibliography}
\end{document}
